% !TEX program = pdflatex
\documentclass[11pt]{article}
\setlength{\textwidth}{6.5in}
\setlength{\oddsidemargin}{0pt}
\setlength{\evensidemargin}{0pt}
\setlength{\textheight}{9in}
\setlength{\topmargin}{0pt}
\addtolength{\topmargin}{-\headheight}
\addtolength{\topmargin}{-\headsep}
\ifdefined\pdfpagewidth \pdfpagewidth=8.5in \pdfpageheight=11in \fi
\ifdefined\pagewidth \pagewidth=8.5in \pageheight=11in \fi
\usepackage{amsmath,amssymb}
\usepackage{array}
\usepackage{booktabs}
\usepackage[tableposition=top]{caption}
\usepackage{longtable}
\usepackage{enumitem}
\usepackage{tikz}
\usepackage{pgfplots}
\pgfplotsset{compat=1.18}
\emergencystretch=3em

% ---- notation ---------------------------------------------------------------
\newcommand{\um}{\,$\mu$m}
\newcommand{\kms}{km\,s$^{-1}$}
\newcommand{\flux}{erg\,cm$^{-2}$\,s$^{-1}$}
\newcommand{\mA}{m\AA}
\newcommand{\sci}[2]{#1\times10^{#2}}
\newcommand{\ion}[2]{#1\,\textsc{#2}}
\newcommand{\nuc}[2]{$^{#1}$#2}
\newcommand{\lxl}{$L_X/L_\mathrm{bol}$}

% ---- front matter -----------------------------------------------------------
\title{Missing Light in Detached Eclipsing Binaries:\\ the Case of KIC 2557430}
\author{Charles C. Norton}
\date{}

\begin{document}
\maketitle

\begin{abstract}
A detached eclipsing binary's light curve and radial velocities give each component's
radius and temperature, and so its luminosity, independently of its distance; with a
Gaia parallax the system's apparent brightness follows, and a structure that intercepts
part of the light would leave the system fainter than its components predict. We test
the $117$ systems of DEBCat with components below $8000$\,K within $1$\,kpc and
renormalised unit weight error (RUWE) of $1.4$ or less. The observed minus predicted
$G$ has a median of $-0.04$\,mag and a robust scatter of $0.12$\,mag; $12$ systems are
too bright, by the light of third stars and blends, and one is too faint by $0.3$\,mag at
$3\sigma$: KIC~2557430, by $1.05\pm0.19$\,mag. Gaia resolves it into the eclipsing pair,
$G=11.44$ at $343$\,pc, and a co-moving K star $2.4''$ away at $339$\,pc, inside
Kepler's pixels. The pair's colour and Gaia temperature, $6130$\,K with no extinction,
lie $900$\,K below the published $7050$\,K primary, and Gaia's luminosities for the two
sources sum to $2.9\,L_\odot$ against the solution's $9.7\,L_\odot$. The light is not
missing from the system: the published temperatures, and the extinction the test inferred
from them, are too high, and a spectroscopic orbit free of the companion's light decides
the radii.
\end{abstract}

% =============================================================================
\section{Introduction}\label{sec:intro}
% =============================================================================

A structure that intercepts part of a star's light and re-radiates it in the infrared
leaves the star fainter in the optical than its size and temperature imply. For a single
star the size comes from models or interferometry, but a detached eclipsing binary
gives both components' radii from its light curve and its velocities, and their
temperatures from spectra and the eclipse depths, with uncertainties of one or two per
cent for the best systems~\cite{Southworth2015}. With a Gaia parallax~\cite{Gaia2023}
the system's luminosity then predicts its apparent magnitude, and a shortfall is light
that the system does not send us. Detached eclipsing binaries have been used the other
way, to test Gaia's parallaxes against their own distances~\cite{Stassun2021}; here
the parallaxes are taken as given and the light is tested.

\subsection{Contributions}
\begin{enumerate}[leftmargin=*,itemsep=1pt]
\item A comparison of the predicted and observed Gaia $G$ magnitudes of the $117$ clean
detached eclipsing binaries of DEBCat.
\item The resolution of the one system too faint, KIC~2557430, by Gaia's view of the
system and its companion.
\end{enumerate}

% =============================================================================
\section{Data and method}\label{sec:method}
% =============================================================================

DEBCat~\cite{Southworth2015} lists $388$ detached eclipsing binaries with masses and
radii measured to a few per cent, and for $252$ of them both components' luminosities
and temperatures and a Gaia DR3 counterpart through CDS XMatch on their SIMBAD names.
For each component, $\log L$ and $\log T_\mathrm{eff}$ give an absolute $G$ through
bolometric corrections in temperature, and the two components summed and carried to
the Gaia distance predict $G$. The extinction applied is the larger of Gaia's $A_G$ and
$2.7\,E(B-V)$, with $E(B-V)$ the catalogue's $B-V$ less the intrinsic colour of the
components' combined light. We keep the $117$ systems whose components are cooler than
$8000$\,K, which lie within $1$\,kpc and whose RUWE is $1.4$ or less, so that the
parallax and the photometry belong to an unresolved, well-measured source; BG~Ind,
IM~Per and the other systems of high RUWE are set aside. A system is too faint when
$G_\mathrm{obs}-G_\mathrm{pred}$ exceeds $0.3$\,mag at $3\sigma$, with the error from
the luminosities, the parallax and the photometry.

% =============================================================================
\section{Results}\label{sec:results}
% =============================================================================

Across the $117$ systems $G_\mathrm{obs}-G_\mathrm{pred}$ has a median of $-0.04$\,mag
and a robust scatter of $0.12$\,mag. Twelve systems are too bright, as light from a
third star or a blend in Gaia's photometry makes them, and one is too faint:
KIC~2557430 (KOI-6277), an F2 and F9 pair with a period of $1.298$\,d, by
$1.05\pm0.19$\,mag ($5.6\sigma$), a factor of $2.6$ in flux.

Its published solution~\cite{Hoyman2020} gives radii of $1.88$ and $1.23\,R_\odot$ and
temperatures of $7050$ and $6050$\,K, luminosities of $7.9$ and $1.8\,L_\odot$, and the
catalogue's $B-V$ of $0.58$ against the $0.38$ of those components set the extinction
in the test at $0.54$\,mag; without it the system would be $1.6$\,mag too faint.
Table~\ref{tab:gaia} gives Gaia's view. The system is the visual double
WDS~J19074+3749AB, and Gaia resolves it into two sources at the same parallax: the
eclipsing pair, $G=11.44$, $G_\mathrm{BP}-G_\mathrm{RP}=0.68$, $2.916\pm0.012$\,mas,
with RUWE $0.81$ and a radial velocity that varies by $20$\,km\,s$^{-1}$ over $20$
transits, and a K star $2.4''$ away, $G=13.30$, $G_\mathrm{BP}-G_\mathrm{RP}=1.39$,
$2.946\pm0.015$\,mas. The test matched the eclipsing pair's source; adding the K star
brings the system to $G=11.26$, still $0.86$\,mag below the prediction.

\begin{table}[t]
\centering\small
\caption{Gaia DR3 for the two sources of KIC~2557430: parallax, $G$, colour, RUWE, and the
temperature, radius and luminosity of the astrophysical parameters (GSP-Phot and FLAME),
which treat each source as a single star.}
\label{tab:gaia}
\begin{tabular}{lrrrrrrr}
\toprule
Source & $\varpi$ (mas) & $G$ & BP$-$RP & RUWE & $T_\mathrm{eff}$ (K) & $R$ ($R_\odot$) & $L$ ($L_\odot$) \\
\midrule
2099133865420292992 & $2.916\pm0.012$ & $11.44$ & $0.68$ & $0.81$ & $6131$ & $1.31$ & $2.20$ \\
2099133865420292480 & $2.946\pm0.015$ & $13.30$ & $1.39$ & $1.55$ & $4698$ & $1.31$ & $0.74$ \\
\bottomrule
\end{tabular}
\end{table}

Gaia's colours tell against the published temperatures rather than for missing light.
Gaia gives the eclipsing pair $A_G=0.001$\,mag, so the catalogue's $B-V$ of $0.58$ is
the pair's own colour, that of a star near $6000$\,K, where a $7050$\,K primary carrying
four fifths of the light would give about $0.4$; Gaia's photometric temperature for the
pair is $6130$\,K. At its parallax the pair's light corresponds to $2.2\,L_\odot$, and
the K star's to $0.7\,L_\odot$, together $2.9\,L_\odot$ against the solution's
$9.7\,L_\odot$. A primary at $6100$\,K with the published radius would hold $4.4\,L_\odot$,
which halves the deficit, and the rest lies within what radii from a spectroscopic orbit
contaminated by a third star can carry, since the K star, $1.9$\,mag fainter in $G$,
falls inside Kepler's $4''$ pixels and any spectrograph's fibre or slit at that
separation. An earlier analysis of the Kepler light curve read the system as a triple
holding a $\gamma$~Doradus pulsator and two flaring components~\cite{Kamil2017}, and the
K star is a flaring component that such a reading needs.

% =============================================================================
\section{Related work}\label{sec:related}
% =============================================================================

Detached eclipsing binaries serve as distance standards, and their luminosities have
been used to test Gaia's parallaxes~\cite{Stassun2021}; searches for intercepted light
around single stars use their infrared excess against models of their
photospheres~\cite{Suazo2024}. We find no earlier search of eclipsing binaries for
light that their measured components should send and do not.

% =============================================================================
\section{Conclusion}\label{sec:conclusion}
% =============================================================================

No clean detached eclipsing binary of DEBCat is fainter than its components predict
beyond the scatter of the comparison, $0.12$\,mag, save KIC~2557430, and Gaia shows
that its deficit belongs to its published solution: the pair's colour and temperature
are those of a $6100$\,K star with no extinction, and a resolved co-moving K star sits
inside the apertures that gave the solution. The test bounds light intercepted from the
other $116$ systems at about $0.36$\,mag, three times the scatter, a third of the flux.
A new spectroscopic orbit of KIC~2557430 that separates the K star's light, with
temperatures from high-resolution spectra, would place it among them.

\section*{Acknowledgements}

This work uses DEBCat, maintained by J.~Southworth; data from the European Space Agency
mission Gaia, processed by the Gaia Data Processing and Analysis Consortium (DPAC), funded
by national institutions, in particular those participating in the Gaia Multilateral
Agreement; the SIMBAD and VizieR services and the XMatch service at CDS, Strasbourg; and
the Washington Double Star Catalog maintained at the U.S. Naval Observatory.

\begin{thebibliography}{99}\small\raggedright
\setlength{\itemsep}{1pt}

\bibitem{Southworth2015}
J.~Southworth, DEBCat: a catalog of detached eclipsing binary stars, in \emph{Living
Together: Planets, Host Stars and Binaries}, ASP Conf. Ser. \textbf{496}, 164 (2015).

\bibitem{Gaia2023}
Gaia Collaboration, A.~Vallenari, A.~G.~A. Brown, et al., Gaia Data Release 3: summary
of the content and survey properties, \emph{Astron. Astrophys.} \textbf{674}, A1 (2023).

\bibitem{Stassun2021}
K.~G. Stassun and G.~Torres, Parallax systematics and photocenter motions of benchmark
eclipsing binaries in Gaia EDR3, \emph{Astrophys. J. Lett.} \textbf{907}, L33 (2021).

\bibitem{Hoyman2020}
B.~Hoyman, \"O.~\c{C}ak\i rl\i\ and O.~\"Ozdarcan, Updated pulsation period relations of
$\gamma$~Dor-type stars in eclipsing binaries: the case of KIC~2557430, \emph{Mon. Not.
R. Astron. Soc.} \textbf{491}, 5980 (2020).

\bibitem{Kamil2017}
C.~Kamil and H.~A. Dal, KIC~2557430: a triple system containing one $\gamma$~Dor and two
flaring components?, \emph{Publ. Astron. Soc. Aust.} \textbf{34} (2017),
arXiv:1707.04727.

\bibitem{Suazo2024}
M.~Suazo, E.~Zackrisson, P.~K. Mahto, et al., Project Hephaistos -- II. Dyson sphere
candidates from Gaia DR3, 2MASS, and WISE, \emph{Mon. Not. R. Astron. Soc.}
\textbf{531}, 695 (2024).

\end{thebibliography}

\end{document}
